\section{A query keeps its domain, endpoints and memory}
\label{sec:query-resolution-boundaries}

The same displayed amount can arise from different input domains, and the
same ending cash can arise from different histories. A query therefore carries
its domain, operation order, endpoint status and requested state. The examples
below make those fields concrete before the implementation is asked to reuse a
summary.

\subsection*{Real amounts, rounded outputs and a cent grid are different models}

Return to the debt family in which nonnegative fixed-rate principal \(F\) and
floating-rate principal \(V\) satisfy
\[
 99.5\le F+V<100.5,\qquad
 4.5\le .05(F+V)<5.5.
\]
The future interest query is \(I=.05F+.07V\). The lower endpoint is attained
at \(F=99.5,V=0\), giving \(I=4.975\). The upper endpoint is the supremum
\(.07(100.5)=7.035\), but it is excluded because the reported total is
strictly below 100.5. The exact real image is
\([4.975,7.035)\), with an attained lower endpoint and an unattained upper
endpoint. A solver result must carry both facts; writing a closed interval
would change the set. To establish that there are no gaps, hold the total at
99.5 and vary \(V\) from zero to 99.5. Then
\(I=4.975+.02V\) covers \([4.975,6.965]\). The all-floating histories
with totals from 99.5 up to, but excluding, 100.5 cover
\([6.965,7.035)\). Their union is the stated interval. Conversely,
\(.05(F+V)\le I\le .07(F+V)\) excludes values outside those bounds.

Now ask for a display rounded to cents, using nearest rounding with exact
halfway values away from zero. Every displayed value from 4.98 through 7.03
has a witness in the real family. The lower raw endpoint displays as 4.98.
The excluded raw supremum would display as 7.04, but no admissible amount
reaches it. Thus the rounded image has 206 values, 4.98--7.03. This is an
image of the real family, not evidence that the underlying amounts are cents.
For any cent amount \(x\) in the image with \(x\le6.965\), a witness is
\(V=(x-4.975)/.02\) and \(F=99.5-V\). Above 6.965 choose
\(F=0,V=x/.07\). These formulas also show why enumerating cent displays
has not discretised the unknown principal.

A different model restricts \(F\) and \(V\) to integer cents before computing
interest. Its maximum is 7.0343 at \(V=100.49\), and its rounded display
still runs from 4.98 to 7.03. Equal displays do not prove equal latent
families: the real model admits values between 7.0343 and 7.035 that the
cent-grid model excludes. A stored result must therefore retain the model
identity and input quantum.

Signed and tie cases expose the same boundary. Negating the real image gives
\((-7.035,-4.975]\); its rounded outputs are the cents from \(-7.03\)
through \(-4.98\). The lower endpoint remains excluded. The interval
\([0,1)\) can round to 1.00 because 0.995 is an admitted value, even though
1 itself is not. A certified enclosure \([1.0049,1.0051]\) crosses the 1.005
tie and cannot select one posting without more information. Unknown tie rule,
zero denominator, and an unproved nonzero denominator are separate failures.
Rounding each of two 0.005 lines before adding gives 0.02, while adding first
and rounding gives 0.01. Operation order belongs to the query method.

\subsection*{A projection is not a reusable transition state}

For a concrete counterexample, take two separately retained histories in which
the timing of a reinsurance treaty changes when income is recognised. Both
end with cash 88.3. Their initial and first-year equities are 85 and 83 for
the timely treaty, but 70 and 74.75 for the late treaty. These are the supplied
comparison values here; Section~\ref{sec:insurance-connected-groups} derives
them from the separate issued and held contracts. The terminal-cash
projection is therefore equal while the transition states and service results
are different. A query requesting only terminal cash may project both histories
to 88.3; a later query requesting first-year equity, loss recovery or treaty
timing must revalidate the resolution and retain the full group states.

The bank commitment supplies the same warning. Funded gross loan 100 does not
contain the original commitment date, credit baseline or observed draw. The
funded allowance 4 cannot be recomputed from that scalar alone. If the original
records were preserved, refinement can recover the state. If they were never
known, the result is \texttt{DETAIL\_UNAVAILABLE} or measurement unresolved,
not a fabricated staging history. A source revision or changed query creates
a new manifest; it does not silently resume an old result.

\subsection*{Conditional mass remains attached to unresolved outcomes}

For a small authored law, assign prior masses \(1/3,2/3\) to immutable timely
and late histories and likelihoods \(3/4,1\) to a newly admitted observation.
The evidence normaliser is \(11/12\), so posterior masses are \(3/11\) and
\(8/11\). The invalid pooled treatment is not a third law atom: it was rejected
as unsupported accounting treatment. If the late history's execution is
unresolved, its mass remains \(8/11\); reporting the timely result with
probability one silently conditions on successful execution. If both
first-year equities are available under that law, their mean is
\((3/11)83+(8/11)74.75=77\). With only the timely branch evaluated, the
known weighted contribution is \((3/11)83\), not the complete mean.
The same event IDs, measure provenance and unresolved mass must travel with
the output.

These examples do not add a general solver or a cent-based accounting policy.
They specify the distinctions that a compatible implementation must preserve:
input domain, output operation, endpoint attainment, event order, retained
economic memory, conditional mass and failure status. A changed output,
horizon, grouping rule or scenario therefore returns
\texttt{RESOLUTION\_REVALIDATION\_REQUIRED} until sufficiency is checked again.
